\documentclass[CJK]{beamer}
\usepackage{CJKutf8}
\usepackage{beamerthemesplit}
\usetheme{Malmoe}
\useoutertheme[footline=authortitle]{miniframes}
\usepackage{amsmath}
\usepackage{amssymb}
\usepackage{graphicx}
\usepackage{eufrak}
\usepackage{color}
\usepackage{slashed}
\usepackage{simplewick}
\usepackage{tikz}
\usepackage{tcolorbox}
\graphicspath{{../figures/}}
%%figures
\def\lfig#1#2{\includegraphics[width=#1 in]{#2}}
\def\addfig#1#2{\begin{center}\includegraphics[width=#1 in]{#2}\end{center}}
\def\wulian{\includegraphics[width=0.18in]{emoji_wulian.jpg}}
\def\bigwulian{\includegraphics[width=0.35in]{emoji_wulian.jpg}}
\def\bye{\includegraphics[width=0.18in]{emoji_bye.jpg}}
\def\bigbye{\includegraphics[width=0.35in]{emoji_bye.jpg}}
\def\huaixiao{\includegraphics[width=0.18in]{emoji_huaixiao.jpg}}
\def\bighuaixiao{\includegraphics[width=0.35in]{emoji_huaixiao.jpg}}
%% colors
\def\blacktext#1{{\color{black}#1}}
\def\bluetext#1{{\color{blue}#1}}
\def\redtext#1{{\color{red}#1}}
\def\darkbluetext#1{{\color[rgb]{0,0.2,0.6}#1}}
\def\skybluetext#1{{\color[rgb]{0.2,0.7,1.}#1}}
\def\cyantext#1{{\color[rgb]{0.,0.5,0.5}#1}}
\def\greentext#1{{\color[rgb]{0,0.7,0.1}#1}}
\def\darkgray{\color[rgb]{0.2,0.2,0.2}}
\def\lightgray{\color[rgb]{0.6,0.6,0.6}}
\def\gray{\color[rgb]{0.4,0.4,0.4}}
\def\blue{\color{blue}}
\def\red{\color{red}}
\def\green{\color{green}}
\def\darkgreen{\color[rgb]{0,0.4,0.1}}
\def\darkblue{\color[rgb]{0,0.2,0.6}}
\def\skyblue{\color[rgb]{0.2,0.7,1.}}
%%control
\def\be{\begin{equation}}
\def\ee{\nonumber\end{equation}}
\def\bea{\begin{eqnarray}}
\def\eea{\nonumber\end{eqnarray}}
\def\bch{\begin{CJK}{UTF8}{gbsn}}
\def\ech{\end{CJK}}
\def\bitem{\begin{itemize}}
\def\eitem{\end{itemize}}
\def\bcenter{\begin{center}}
\def\ecenter{\end{center}}
\def\bex{\begin{minipage}{0.2\textwidth}\includegraphics[width=0.6in]{jugelizi.png}\end{minipage}\begin{minipage}{0.76\textwidth}}
\def\eex{\end{minipage}}
\def\chtitle#1{\frametitle{\bch#1\ech}}
\def\bmat#1{\left(\begin{array}{#1}}
\def\emat{\end{array}\right)}
\def\bcase#1{\left\{\begin{array}{#1}}
\def\ecase{\end{array}\right.}
\def\bmini#1{\begin{minipage}{#1\textwidth}}
\def\emini{\end{minipage}}
\def\tbox#1{\begin{tcolorbox}#1\end{tcolorbox}}
\def\pfrac#1#2#3{\left(\frac{\partial #1}{\partial #2}\right)_{#3}}
%%symbols
\def\sone{$\star$}
\def\stwo{$\star\star$}
\def\sthree{$\star\star\star$}
\def\sfour{$\star\star\star\star$}
\def\sfive{$\star\star\star\star\star$}
\def\rint{{\int_\leftrightarrow}}
\def\roint{{\oint_\leftrightarrow}}
\def\stdHf{{\textit{\r H}_f}}
\def\deltaH{{\Delta \textit{\r H}}}
\def\ii{{\dot{\imath}}}
\def\skipline{{\vskip0.1in}}
\def\skiplines{{\vskip0.2in}}
\def\lagr{{\mathcal{L}}}
\def\hamil{{\mathcal{H}}}
\def\vecv{{\mathbf{v}}}
\def\vecx{{\mathbf{x}}}
\def\vecy{{\mathbf{y}}}
\def\veck{{\mathbf{k}}}
\def\vecp{{\mathbf{p}}}
\def\vecn{{\mathbf{n}}}
\def\vecA{{\mathbf{A}}}
\def\vecP{{\mathbf{P}}}
\def\vecsigma{{\mathbf{\sigma}}}
\def\hatJn{{\hat{J_\vecn}}}
\def\hatJx{{\hat{J_x}}}
\def\hatJy{{\hat{J_y}}}
\def\hatJz{{\hat{J_z}}}
\def\hatj#1{\hat{J_{#1}}}
\def\hatphi{{\hat{\phi}}}
\def\hatq{{\hat{q}}}
\def\hatpi{{\hat{\pi}}}
\def\vel{\upsilon}
\def\Dint{{\mathcal{D}}}
\def\adag{{\hat{a}^\dagger}}
\def\bdag{{\hat{b}^\dagger}}
\def\cdag{{\hat{c}^\dagger}}
\def\ddag{{\hat{d}^\dagger}}
\def\hata{{\hat{a}}}
\def\hatb{{\hat{b}}}
\def\hatc{{\hat{c}}}
\def\hatd{{\hat{d}}}
\def\hatN{{\hat{N}}}
\def\hatH{{\hat{H}}}
\def\hatp{{\hat{p}}}
\def\Fup{{F^{\mu\nu}}}
\def\Fdown{{F_{\mu\nu}}}
\def\newl{\nonumber \\}
\def\vece{\mathrm{e}}
\def\calM{{\mathcal{M}}}
\def\calT{{\mathcal{T}}}
\def\calR{{\mathcal{R}}}
\def\barpsi{\bar{\psi}}
\def\baru{\bar{u}}
\def\barv{\bar{\upsilon}}
\def\qeq{\stackrel{?}{=}}
\def\torder#1{\mathcal{T}\left(#1\right)}
\def\rorder#1{\mathcal{R}\left(#1\right)}
\def\contr#1#2{\contraction{}{#1}{}{#2}#1#2}
\def\trof#1{\mathrm{Tr}\left(#1\right)}
\def\trace{\mathrm{Tr}}
\def\comm#1{\ \ \ \left(\mathrm{used}\ #1\right)}
\def\tcomm#1{\ \ \ (\text{#1})}
\def\slp{\slashed{p}}
\def\slk{\slashed{k}}
\def\calp{{\mathfrak{p}}}
\def\veccalp{\mathbf{\mathfrak{p}}}
\def\Tthree{T_{\tiny \textcircled{3}}}
\def\pthree{p_{\tiny \textcircled{3}}}
\def\dbar{{\,\mathchar'26\mkern-12mu d}}
\def\erf{\mathrm{erf}}
\def\const{\mathrm{constant}}
\def\pheat{\pfrac p{\ln T}V}
\def\vheat{\pfrac V{\ln T}p}
%%units
\def\fdeg{{^\circ \mathrm{F}}}
\def\cdeg{^\circ \mathrm{C}}
\def\atm{\,\mathrm{atm}}
\def\angstrom{\,\text{\AA}}
\def\SIL{\,\mathrm{L}}
\def\SIkm{\,\mathrm{km}}
\def\SIyr{\,\mathrm{yr}}
\def\SIGyr{\,\mathrm{Gyr}}
\def\SIV{\,\mathrm{V}}
\def\SImV{\,\mathrm{mV}}
\def\SIeV{\,\mathrm{eV}}
\def\SIkeV{\,\mathrm{keV}}
\def\SIMeV{\,\mathrm{MeV}}
\def\SIGeV{\,\mathrm{GeV}}
\def\SIcal{\,\mathrm{cal}}
\def\SIkcal{\,\mathrm{kcal}}
\def\SImol{\,\mathrm{mol}}
\def\SIN{\,\mathrm{N}}
\def\SIHz{\,\mathrm{Hz}}
\def\SIm{\,\mathrm{m}}
\def\SIcm{\,\mathrm{cm}}
\def\SIfm{\,\mathrm{fm}}
\def\SImm{\,\mathrm{mm}}
\def\SInm{\,\mathrm{nm}}
\def\SImum{\,\mathrm{\mu m}}
\def\SIJ{\,\mathrm{J}}
\def\SIW{\,\mathrm{W}}
\def\SIkJ{\,\mathrm{kJ}}
\def\SIs{\,\mathrm{s}}
\def\SIkg{\,\mathrm{kg}}
\def\SIg{\,\mathrm{g}}
\def\SIK{\,\mathrm{K}}
\def\SImmHg{\,\mathrm{mmHg}}
\def\SIPa{\,\mathrm{Pa}}

\def\courseurl{https://github.com/zqhuang/SYSU\_TD}

\def\tpage#1#2{
\begin{frame}
\begin{center}
\begin{Large}
\bch
热学 \\
第#1讲 #2

{\vskip 0.3in}

黄志琦

\ech
\end{Large}
\end{center}

\vskip 0.2in

\bch
教材：《热学》第二版，赵凯华，罗蔚茵，高等教育出版社
\ech

\bch
课件下载
\ech
\courseurl
\end{frame}
}

\title{Lesson 13 Conditions for Equilibrium }
  \author{}
  \date{}
\begin{document}
\tpage{13}{平衡判据}

\section{Review}

\begin{frame}
\chtitle{上讲内容回顾}
\bch 
\bitem
\item{热力学第二定律的四种表述}
\item{克劳修斯不等式和热力学第二定律的“微分”表示}
\item{熵的计算}
\eitem
\ech
\end{frame}

\begin{frame}
\chtitle{本讲内容}
\bch
\bitem
\item{眼花撩乱的平衡判据和背书宝的福音}
\item{真的懂了吗}
\eitem
\ech
\end{frame}


\section{A Smart Rule}

\begin{frame}
\chtitle{眼花撩乱的平衡判据}
\bch
教材从198页开始叨念各种（完全看不懂的）“平衡判据”：

\bitem
\item{熵判据：在内能和体积不变的情况下，对于一切可能的变动来说，平衡态的熵最大。}
\item{自由能判据：在温度和体积不变的情况下，对于一切可能的变动来说，平衡态的自由能最小。}
\item{自由焓判据：在温度和压强不变的情况下，对于一切可能的变动来说，平衡态的自由焓最小。}
\eitem

\addfig{1.5}{bzdsh.jpg}
\ech
\end{frame}

\begin{frame}
\chtitle{背书宝的福音}
\bch
\addfig{1}{beishubao.jpg}

要写出“$X$判据”，

\bitem
\item{第一步：把$X$写成全微分：$dX = \ldots dY + \ldots dZ$}
\item{第二步：背书。在$Y$和$Z$不变的情况下，对于一切可能的变动来说，平衡态的$X$最大/小。}
\item{根据常识选择“最大”还是“最小”。熵当然取最大，各种能量和焓当然取最小。}
\eitem

似乎有点理解了：$Y$和$Z$不变，则$dX = 0$，代表$X$是某种极值？
\ech
\end{frame}


\begin{frame}
\chtitle{例子1：熵判据}
\bch
$$dS = \frac{1}{T} dU + \frac{p}{T} dV $$
背书：

在内能和体积不变的情况下，对于一切可能的变动来说，平衡态的熵最大。

\skiplines

熵判据也常常说成，孤立系的平衡态熵最大。内能和体积不变，是孤立系的一种。
\ech
\end{frame}


\begin{frame}
\chtitle{例子2：自由能判据}
\bch
$$dF = -S dT - p dV $$
背书：

在温度和体积（包括形状）不变的情况下，对于一切可能的变动来说，平衡态的自由能最小。

\ech
\end{frame}


\begin{frame}
\chtitle{例子3：自由焓判据}
\bch
$$dG = -S dT + V dp$$
背书：

在温度和压强不变的情况下，对于一切可能的变动来说，平衡态的自由焓最小。

\ech
\end{frame}


\begin{frame}
\chtitle{为什么没见过“内能判据”和“焓判据”？}
\bch
$$dU = TdS - pdV$$
$$dH = TdS + Vdp $$

如果要背书：在熵和$\ldots$不变的情况下,$\ldots$

\skipline

等等，怎么扰动系统而又保持熵不变？\wulian这个实验上操作起来比较难，所以这样的判据虽然存在但并不常被提到。
\ech
\end{frame}


\section{Understanding the Statements}

\begin{frame}
\chtitle{真懂了吗}
\bch
“一切可能的变动”是什么意思？

\addfig{1.5}{zhendongma.jpg}
\ech
\end{frame}

\begin{frame}
\chtitle{理解熵判据}
\bch
熵判据里“一切可能的变动”好理解：就是指孤立系的平衡态和一切非平衡态。对非平衡态，广延量内能和体积都有定义。


\addfig{1}{meimaobing.jpg}

证明：根据孤立系的熵增大原理，任何非平衡态的熵会持续增大直至到达平衡态。显然平衡态的熵是最大的。

\ech
\end{frame}


\begin{frame}
\chtitle{理解自由能判据}
\bch
自由能判据的条件是定温定体。“一切可能的变动”也包括非平衡态。这时需要假设系统和恒温的理想传导热库接触来理解非平衡态的温度。

\addfig{1}{meimaobing.jpg}

证明：从非平衡态到平衡态的过程为不可逆过程。定体过程不对外做功，又由于恒温条件下对外做功小于自由能的减少（对不可逆过程）。故自由能的减少大于零。即非平衡态的自由能大于平衡态的自由能。

\ech
\end{frame}

\begin{frame}
\chtitle{理解自由焓判据}
\bch
自由焓判据的条件是定温定压。“一切可能的变动”也包括非平衡态。这时需要假设系统和恒温的理想传导热库接触来理解非平衡态的温度。

\addfig{1}{meimaobing.jpg}

证明：从非平衡态到平衡态的过程为不可逆过程。定温定压条件下，对不可逆过程，其他形式的做功小于自由焓的减少，又这里不考虑其他形式的做功（即没有额外安装发电机什么的），故自由焓的减少大于零。即非平衡态的自由焓大于平衡态的自由焓。

\ech
\end{frame}

\begin{frame}
\chtitle{三大平衡条件}
\bch
\bitem
\item[1]{\blue 热平衡条件：系统内部温度均匀}
\item[2]{\blue 力学平衡条件：系统内部压强均匀}
\item[3]{\blue 相平衡条件：系统内各相化学势相等}
\eitem

{\scriptsize
注：第三条仅考虑了物态变化（如液态水和水蒸气之间的平衡），其实它可以推广到涉及化学反应的一般的化学平衡条件。化学反应比相变更为复杂，各种反应物或者产物有可能混合在一起产生额外的混合熵（请参考教材208页）。
}
\ech
\end{frame}


\begin{frame}
\chtitle{热平衡条件}
\bch
热平衡条件可以从孤立系的熵判据推出：如果有相邻的温度不相等的子系统，则由高温到低温的热量传输必然增大熵。这说明系统还未达到平衡态（熵最大的状态）。
\addfig{1.5}{eqc1.jpg}
$$\Delta S = \frac{-Q}{T_1} + \frac{Q}{T_2} > 0$$
\ech
\end{frame}

\begin{frame}
\chtitle{力学平衡条件}
\bch
力学平衡条件可以从定温定体系统的自由能判据推出：如果有相邻的压强不相等的子系统，则压强大的子系统等温膨胀$\Delta V$，压强小的子系统等温收缩$\Delta V$，可使系统温度和体积不变的情况下总自由能减小。这说明系统还未达到平衡态（自由能最低的态）。
\addfig{1.5}{eqc2.jpg}
$$\Delta F = -p_1 \Delta V + p_2 \Delta V < 0 $$
\ech
\end{frame}


\begin{frame}
\chtitle{相平衡条件}
\bch
设有种可以相互转化的物态。每种物态的粒子数分别为$N_1$, $N_2$, $\ldots$, $N_n$。每种物态成分的化学势即该成分的平均每个粒子的自由焓，即
$$ \mu_i \equiv \pfrac G{N_i}{T,p,N_1, N_2, \ldots, N_{i-1}, N_{i+1}, \ldots, N_n} $$

\skipline

相平衡条件可以从定温定压系统的自由焓判据推出：如果定温定压条件下$ \mu_i > \mu_j$，则可以通过$i\rightarrow j$的相变降低系统的总自由焓。
，这说明系统还未达到平衡态。
\ech
\end{frame}


\begin{frame}
\chtitle{经典理想气体的化学势}
\bch
{\small
考虑经典的理想气体，在$i$态上的分子数为
$$ n_i = e^{-\frac{\varepsilon_i-\mu}{kT}}$$
下面我们证明，这里的$\mu$和用自由焓定义的化学势相同}

{\scriptsize
设分子数为$N$。在固定压强和温度的情况下，由$pV = NkT$知体积
$V \propto N$， 且
$$\pfrac VN{p,T} = \frac{kT}{p} $$

由理想气体的熵变公式
$$ \Delta S = \int \frac{C_V}{T} dT + \nu R \Delta \ln V $$
由于$\Delta \ln V =\Delta \ln N$，$\nu R = Nk$，上式就可以写成
$$ S = S_0 + N k \ln \frac{N}{N_0} $$
故
$$\pfrac SN{p,T} = k (\ln \frac{N}{N_0} + 1) $$




}
\ech
\end{frame}

\end{document}
